\section*{Chapitre \ref{ch:b2:atmospheric-thermodynamics} --- Thermodynamique de l'atmosphère}

\begin{solution}{exo:b2:atmospheric-thermodynamics:1}
(a) $e = e_s(\qty{278.15}{K}) = \qty{8.73}{hPa}$ et $e_s(\qty{283.15}{K}) = \qty{12.32}{hPa}$.
(b) $8.73/12.32 = \qty{70.9}{\percent}$. (c) $r = 0.622 \times 8.73/(850 - 8.73) =
\qty{6.45}{g/kg}$ et $q = r/(1 + r) = \qty{6.41}{g/kg}$.
\end{solution}

\begin{solution}{exo:b2:atmospheric-thermodynamics:2}
(a) $\mathrm{d}\ln e_s/\mathrm{d}T = 5419/T^2$~: \qty{7.26}{\percent} par kelvin à
\qty{273.15}{K}, \qty{6.31}{\percent} à \qty{293.15}{K}, \qty{5.90}{\percent} à \qty{303.15}{K}.
(b) $e_s(\qty{20}{\celsius}) = \qty{23.66}{hPa}$ et $e_s(\qty{25}{\celsius})
= \qty{32.26}{hPa}$~; rapport 1.363, contre $1.063^5 = 1.357$. Le taux est une dérivée
de $\ln e_s$, donc sur \qty{5}{K} le facteur vaut $\exp(5 \times 0.0620) = 1.363$, avec
0.0620 le taux au point milieu, \qty{22.5}{\celsius}~; capitaliser \qty{6.3}{\percent}
cinq fois donne un résultat proche mais non exact. (c) $23.66/23.39 = 1.012$~: la
formule est trop forte de \qty{1.2}{\percent}, car elle garde la chaleur latente du point
triple, alors que la chaleur latente décroît quand la température s'élève~; la pente
de la formule est trop raide, et son erreur croît avec la distance à \qty{0}{\celsius}.
\end{solution}

\begin{solution}{exo:b2:atmospheric-thermodynamics:3}
(a) $125 \times (30 - 14) = \qty{2000}{m}$. (b) $30 - 9.8 \times 2.0 = \qty{10.4}{\celsius}$. (c)
$14 - 1.8 \times 2.0 = \qty{10.4}{\celsius}$~: la température et le point de rosée se rejoignent,
comme ils le doivent au niveau de saturation~; la règle est ce calcul. (d)
$125 \times (33 - 14) = \qty{2375}{m}$~: les bases montent à mesure que l'après-midi réchauffe l'air de surface.
\end{solution}

\begin{solution}{exo:b2:atmospheric-thermodynamics:4}
(a) La courbe du point de rosée se confond avec celle de la température depuis le sol jusqu'à environ
\qty{908}{hPa}~: la tranche la plus basse, d'un demi-kilomètre, est saturée, ou presque. (b) Entre environ
908 et \qty{889}{hPa}, la courbe de température tourne vers la droite, vers des valeurs plus chaudes,
à mesure que la pression baisse~: la température croît avec l'altitude. (c) LCL vers \qty{964}{hPa}, à
quelques dizaines de mètres au-dessus du sol~; LFC vers \qty{730}{hPa}~; EL vers \qty{235}{hPa}. (d) Au
sommet de l'inversion, le point de rosée est \qty{11}{K} sous la température, et \qty{20}{K} sous elle
à \qty{850}{hPa}~: de l'air chaud et très sec repose sur de l'air frais et humide. Son humidité relative
est faible, donc il n'a pas acquis sa chaleur par condensation~; l'air qui descend se réchauffe par
compression au taux adiabatique sec en gardant son \omterm{def:b2:atmospheric-thermodynamics:moist-air}{rapport de mélange}, ce qui donnerait une telle couche.
\end{solution}

\begin{solution}{exo:b2:atmospheric-thermodynamics:5}
(a) $T_v = 303.15 \times (1 + 0.020/0.622)/1.020 = \qty{306.76}{K}$~; l'approximation donne
$303.15 \times (1 + 0.61 \times 0.020) = \qty{306.85}{K}$. (b) $\rho = p/(R_dT_v) = 10^5/(287.0
\times 306.76) = \qty{1.1358}{kg/m^3}$, contre $10^5/(287.0 \times 303.15) =
\qty{1.1494}{kg/m^3}$ pour l'air sec : \qty{1.2}{\percent} plus léger. (c) Par $T_v - T =
\qty{3.6}{K}$. (d) Une particule qui monte n'a besoin que d'être plus légère que son
environnement~: au-dessus des mers chaudes, l'humidité de l'air ajoute une
flottabilité équivalant à plusieurs kelvins, autant qu'un excès de température notable.
\end{solution}

\begin{solution}{exo:b2:atmospheric-thermodynamics:6}
(a) $e_s$ = 6.11, 23.66 et \qty{43.54}{hPa}~; $r_s = 0.622\,e_s/(1000 - e_s)$ = 3.82,
15.07 et \qty{28.31}{g/kg}~; $\Gamma_s$ = 6.48, 4.20 et \qty{3.44}{K/km}. (b) La particule
sèche se refroidit de \qty{9.8}{K}, la particule saturée de \qty{4.2}{K}~: la particule
saturée est \qty{5.6}{K} plus chaude. (c) Plus haut, l'air est plus froid, $r_s$ est
plus petit, il reste peu de vapeur à condenser, et $\Gamma_s$ revient vers $\Gamma_d$.
\end{solution}

\begin{solution}{exo:b2:atmospheric-thermodynamics:7}
(a) $e = e_s(\qty{291.65}{K}) = \qty{21.52}{hPa}$, $r = 0.622 \times 21.52/(971.1 - 21.52) =
\qty{14.1}{g/kg}$~; $\theta = 292.15 \times (1000/971.1)^{0.2857} = \qty{294.6}{K}$~; $\theta_e =
294.6 \exp[2.501 \times 10^6 \times 0.0141/(1004.5 \times 291.5)] = 294.6 \times 1.128 =
\qty{332.3}{K}$. (b) $e_s(\qty{264.25}{K}) = \qty{3.13}{hPa}$, $r_s = 0.622 \times 3.13/(500 -
3.13) = \qty{3.92}{g/kg}$~; $\theta = 264.25 \times 2^{0.2857} = \qty{322.1}{K}$~; la particule
est saturée, donc $T_{\text{LCL}} = T$ et $\theta_e = 322.1 \exp[2.501 \times 10^6 \times
0.00392/(1004.5 \times 264.25)] = \qty{334.3}{K}$. (c) $\theta$ a augmenté de \qty{27.5}{K},
réchauffé par la chaleur latente des \qty{10.2}{g/kg} de vapeur condensée~; $\theta_e$ a
varié de \qty{0.6}{\percent}, dans la précision de sa formule approchée~: elle est conservée.
\end{solution}

\begin{solution}{exo:b2:atmospheric-thermodynamics:8}
(a) $\text{CAPE} = 287.0 \times 4 \times \ln(700/300) = 287.0 \times 4 \times 0.847 =
\qty{973}{J/kg}$. (b) $\sqrt{2 \times 973} = \qty{44.1}{m/s}$. (c) Entre 1 et \qty{1500}{J/kg}~:
positive, mais pas grande. (d) Une bande de largeur horizontale constante entre deux courbes~;
d'après la \cref{prop:b2:atmospheric-thermodynamics:areas}, son aire est proportionnelle à $\int
\Delta T\,\mathrm{d}\ln p = 4 \ln(700/300)$, quelle que soit la pente des isothermes.
\end{solution}

\begin{solution}{exo:b2:atmospheric-thermodynamics:9}
(a) $D_c = \sqrt{3 \times 0.61 \times 80^3/2.1}\,\unit{nm} = \qty{668}{nm} =
\qty{0.668}{\micro m}$~; $s_c = 2A/(3D_c) = 2 \times 2.1/(3 \times 668) = 0.00210 =
\qty{0.210}{\percent}$. (b) En résolvant $s_c = \sqrt{4A^3/(27\kappa D_d^3)}$ pour $D_d$~: $D_d =
\bigl(4A^3/(27\kappa s_c^2)\bigr)^{1/3}$, \qty{63}{nm} à \qty{0.3}{\percent} et \qty{131}{nm}
à \qty{0.1}{\percent}. (c) Toute particule comprise entre 63 et \qty{131}{nm} s'active alors
aussi~: davantage de gouttelettes et, pour la même eau liquide, des gouttelettes plus petites.
\end{solution}

\begin{solution}{exo:b2:atmospheric-thermodynamics:10}
(a) $w = \sqrt{2 \times 130} = \qty{16.1}{m/s}$, impossible pour un air au repos près du sol. (b)
La température potentielle au sommet de l'inversion vaut $293.65 \times (1000/889.1)^{0.2857} =
\qty{303.7}{K}$~; une particule de surface de même $\theta$ a $T = 303.7 \times (971.1/1000)^{0.2857}
= \qty{301.2}{K}$, soit \qty{28.0}{\celsius}. (c) $125 \times (28.0 - 18.5) = \qty{1190}{m}$ au-dessus
du sol, au-dessus de l'ancien sommet de l'inversion. (d) Un front, qui soulève l'air sur sa pente~; la
lente ascension de l'air situé devant une perturbation en altitude, qui refroidit le couvercle lui-même.
\end{solution}

\begin{solution}{exo:b2:atmospheric-thermodynamics:11}
(a) Lorsque $r_s \to 0$, le numérateur tend vers $g$ et le dénominateur vers $c_p$. (b) Écrire
$L^2r_s/(R_vT^2) = c_p \times \bigl(Lr_s/(R_dT)\bigr) \times LR_d/(R_vTc_p)$, avec
$LR_d/(R_vTc_p) = 5.1$ à \qty{303}{K}~: lorsque $Lr_s/(R_dT) \gg 1$, aussi $L^2r_s/(R_vT^2) \gg c_p$,
et $\Gamma_s \to g\,(Lr_s/(R_dT))/(L^2r_s/(R_vT^2)) = gR_vT/(LR_d) = 9.81 \times 461.5 \times
303/(2.501 \times 10^6 \times 287.0) = \qty{1.91}{K/km}$. (c) Avec $T_d = T$, $\Gamma_{\text{rosée}} = gR_vT/(LR_d)$,
le même. Dans cette limite, la chaleur latente est si grande que la particule ne
peut condenser presque aucune de sa vapeur sans se réchauffer bien au-dessus de son point de saturation~:
elle garde son \omterm{def:b2:atmospheric-thermodynamics:moist-air}{rapport de mélange} et suit la droite de \omterm{def:b2:atmospheric-thermodynamics:moist-air}{rapport de mélange} constant, le long de
laquelle la température baisse à $\Gamma_{\text{rosée}}$. (d) À \qty{30}{\celsius} et \qty{1000}{hPa},
$Lr_s/(R_dT) = 0.81$, pas beaucoup plus grand que 1, et $L^2r_s/(R_vT^2) = 4180$, quatre fois
$c_p$~: $\Gamma_s = \qty{3.4}{K/km}$, encore loin de \qty{1.9}{K/km}. La limite n'est pas atteinte.
\end{solution}

\begin{solution}{exo:b2:atmospheric-thermodynamics:12}
(a) $2.868/2.599 = 1.104$ et $1.257/1.032 = 1.218$~: \qty{10.4}{\percent} et \qty{21.8}{\percent}.
(b) Les gouttelettes s'activent à quelques dixièmes de pour cent de \omterm{def:b2:atmospheric-thermodynamics:microphysics}{sursaturation},
des dizaines de fois moins. (c) Les cristaux grandissent vite par dépôt~; ils prennent la
vapeur de l'air, qui passe sous la saturation au-dessus des gouttelettes~; les gouttelettes
s'évaporent et leur vapeur va à la glace (le processus de Bergeron), jusqu'à ce que le nuage
soit entièrement glacé. (d) À \qty{0}{\celsius}, les deux \omterm{def:b2:atmospheric-thermodynamics:moist-air}{pressions de vapeur} saturante
sont égales~: elles se rejoignent au point triple, et le rapport tend vers un près de lui.
\end{solution}

\begin{solution}{pb:b2:atmospheric-thermodynamics:1}
\textbf{1.} $e = e_s(\qty{291.65}{K}) = \qty{21.5}{hPa}$.
\textbf{2.} $r = 0.622 \times 21.5/(971.1 - 21.5) = \qty{14.1}{g/kg}$.
\textbf{3.} $q = 14.1/1.0141 = \qty{13.9}{g/kg}$.
\textbf{4.} $e_s(\qty{292.15}{K}) = \qty{22.2}{hPa}$; $21.5/22.2 = \qty{97}{\percent}$.
\textbf{5.} $0.61 \times 0.0141 \times 292.15 = \qty{2.5}{K}$.
\textbf{6.} $125 \times 0.5 = \qty{63}{m}$; $p = 971.1 \exp[-9.81 \times 62.5/(287.0 \times 294.7)]
= \qty{964}{hPa}$.
\textbf{7.} $\theta = 292.15 \times (1000/971.1)^{0.2857} = \qty{294.6}{K}$.
\textbf{8.} $16.9 - 17.3 = \qty{-0.4}{K}$~: plus froid.
\textbf{9.} $(20.5 - 16.4)/(1102 - 924) = 4.1/178$~: la température monte de \qty{23}{K/km}, soit
un gradient de \qty{-23}{K/km}. Une inversion~: absolument stable, un couvercle.
\textbf{10.} $15.4 - 20.5 = \qty{-5.1}{K}$.
\textbf{11.} $e_s(\qty{-1.8}{\celsius}) = \qty{5.4}{hPa}$ et $e_s(\qty{18.6}{\celsius}) =
\qty{21.7}{hPa}$~: \qty{25}{\percent}. Chaud et sec, il peut s'agir d'air descendu qui s'est réchauffé par
compression, ou d'air chauffé au-dessus d'un sol sec en amont~; ce n'est pas l'air humide situé dessous.
\textbf{12.} $18.6 + 18.4 = \qty{37.0}{K}$ sur $5725 - 1489 = \qty{4236}{m}$, soit \qty{8.73}{K/km}~: entre $\Gamma_s$ et
$\Gamma_d$, en \omterm{def:b2:atmospheric-thermodynamics:instability}{instabilité conditionnelle}.
\textbf{13.} $-18.4 - (-8.9) = \qty{-9.5}{K}$~: instabilité extrême ($-8$ à $-10$).
\textbf{14.} Du sol vers le haut~: 0, $-0.4$, $-0.2$, $-5.1$, $-4.9$, $+1.6$, $+9.5$, $+10.0$,
$+5.6$, $+1.5$, $-9.5$~\unit{K}.
\textbf{15.} Entre 850 ($-4.9$) et \qty{700}{hPa} ($+1.6$), à la fraction $4.9/6.5 = 0.754$~:
$\ln p = \ln 850 + 0.754 \ln(700/850)$, $p = \qty{734}{hPa}$; $z = 1489 + 0.754 \times 1621 =
\qty{2711}{m}$, \qty{2.37}{km} au-dessus du sol.
\textbf{16.} Entre 250 ($+1.5$) et \qty{200}{hPa} ($-9.5$), à la fraction $1.5/11.0 =
0.136$~: $p = \qty{242.5}{hPa}$ et $z = 10560 + 0.136 \times 1446 = \qty{10757}{m}$, environ
\qty{200}{m} sous la tropopause indiquée (\qty{10954}{m}).
\textbf{17.} Les segments (moyenne de $\Delta T$ fois le $\ln$ du rapport des pressions)~: $0.80 \times 0.0474
= 0.038$; $5.55 \times 0.3365 = 1.867$; $9.75 \times 0.2231 = 2.176$; $7.80 \times 0.2877 =
2.244$~; $3.55 \times 0.1823 = 0.647$~; $0.75 \times 0.0305 = 0.023$~; somme \qty{6.995}{K}.
$\text{CAPE} = 287.0 \times 6.995 = \qty{2010}{J/kg}$, à trois chiffres significatifs.
\textbf{18.} Entre 1500 et \qty{2500}{J/kg}~: grande.
\textbf{19.} $-0.2 \times 0.0486 - 0.3 \times 0.0189 - 2.65 \times 0.0207 - 5.0 \times 0.0450 -
2.45 \times 0.1467 = -0.655$~\unit{K}; $\text{CIN} = 287.0 \times (-0.655) = \qty{-188}{J/kg}$.
\textbf{20.} Le tableau ne garde que onze niveaux~: entre eux, la méthode des trapèzes remplace les
courbes par des segments droits en $\ln p$ et en rate les inflexions. La grille de \qty{1}{hPa} suit les courbes. La \omterm{def:b2:atmospheric-thermodynamics:virtual}{température virtuelle}
ajoute la flottabilité de la vapeur d'eau de la particule, plus abondante dans la particule saturée que dans
l'air sec qui l'entoure~: une CAPE plus grande, une CIN plus petite.
\textbf{21.} $\sqrt{2 \times 2010} = \qty{63}{m/s}$.
\textbf{22.} $10757 - 2711 = \qty{8046}{m}$ à $63.4/2 = \qty{31.7}{m/s}$~: \qty{254}{s}, environ quatre
minutes.
\textbf{23.} Le poids de l'eau condensée que transporte le courant ascendant~; le mélange d'air plus sec
qui l'entoure (entraînement), qui le refroidit lorsque les gouttelettes s'évaporent.
\textbf{24.} La théorie de la particule donne au courant ascendant d'un orage alimenté par cet air une vitesse d'environ
\textbf{\qty[text-series-to-math=true]{63}{m.s^{-1}}}, une borne supérieure.
\end{solution}
